\documentclass[10pt]{article}
\usepackage[T1]{fontenc}
\usepackage[utf8]{inputenc}
\usepackage[french]{babel}
\usepackage{amsmath,amssymb,mathrsfs,mathpazo}
\usepackage{geometry}
\geometry{a4paper,margin=2cm}
\usepackage{graphicx,array,tabularx,tabularray,colortbl,multirow,multicol,booktabs,enumitem,pifont,ulem,fancybox,diagbox,cancel}
\usepackage{pstricks,pstricks-add,pst-plot,pst-tree,pst-node,pst-text,pst-func,pst-math,pst-bspline,pst-3dplot,pst-eucl,pst-solides3d}
\let\pscalculate\undefined
\usepackage{pst-bezier}
\usepackage{tikz,pgfplots,tkz-tab}
\usetikzlibrary{arrows,arrows.meta,patterns,calc,positioning,trees,decorations.markings,shapes.geometric,shapes.misc,shapes.arrows,intersections,decorations.pathmorphing,decorations.pathreplacing,backgrounds,shadows,babel}
\pgfplotsset{compat=1.18}
\usepgfplotslibrary{fillbetween}
\usepackage[np]{numprint}
\usepackage{listings,makecell,pas-tableur,twoopt}
\let\barre\undefined
\usepackage[pstricks]{tablvar}
\definecolor{darkgreen}{RGB}{0,100,0}
\providecommand{\red}{\color{red}}
\providecommand{\blue}{\color{blue}}
\providecommand{\green}{\color{green}}
\providecommand{\black}{\color{black}}
\providecommand{\white}{\color{white}}
\providecommand{\gray}{\color{gray}}
\providecommand{\cyan}{\color{cyan}}
\providecommand{\magenta}{\color{magenta}}
\providecommand{\yellow}{\color{yellow}}
\providecommand{\violet}{\color{violet}}
\providecommand{\orange}{\color{orange}}
\providecommand{\euro}{\texteuro}
\providecommand{\EUR}{\texteuro}
\providecommand{\ieme}{\textsuperscript{e}}
\providecommand{\degres}{\ensuremath{^\circ}}
\providecommand{\degree}{\ensuremath{^\circ}}
\providecommand{\vv}[1]{\overrightarrow{#1}}
\providecommand{\Coord}[1]{#1}
\providecommand{\decosix}{}
\providecommand{\starredbullet}{\ensuremath{\bullet}}
\pagestyle{empty}
\setlength{\parindent}{0pt}
\usepackage{adjustbox}\DeclareUnicodeCharacter{2236}{:}\DeclareUnicodeCharacter{2190}{\ensuremath{\leftarrow}}
\providecommand{\R}{\mathbb{R}}
\providecommand{\N}{\mathbb{N}}
\providecommand{\Z}{\mathbb{Z}}
\providecommand{\Q}{\mathbb{Q}}
\providecommand{\D}{\mathbb{D}}
\providecommand{\C}{\mathbb{C}}
\providecommand{\e}{{\mathrm{e}}}
\providecommand{\vect}[1]{\overrightarrow{#1}}
\providecommand{\vectt}[1]{\overrightarrow{\mathrm{#1}}}
\providecommand{\vvt}[1]{\overrightarrow{\mathrm{#1}}}
\providecommand{\np}[2][]{{#2~#1}}
\providecommand{\vv}[1]{{\overrightarrow{#1}}}
\renewcommand{\d}{{\mathrm{d}}}
\providecommand{\degres}{{^\circ}}
\providecommand{\degre}{{^\circ}}
\providecommand{\blue}{\color{blue}}
\providecommand{\red}{\color{red}}
\providecommand{\enskip}{\quad}
\providecommand{\llbracket}{\lbrack\!\lbrack}
\providecommand{\rrbracket}{\rbrack\!\rbrack}
\providecommand{\euro}{{\text{€}}}
\providecommand{\no}{{\text{n°}}}
\providecommand{\interval}[2]{{[#1;#2]}}
\providecommand{\intervFO}[2]{{[#1;#2[}}
\providecommand{\dotfill}{\ldots}

\providecommand{\ds}{\displaystyle}
\providecommand{\textding}[1]{#1}
\providecommand{\Oij}{\(\mathrm{O};\vec{\imath},\vec{\jmath}\)}
\providecommand{\Oijk}{\(\mathrm{O};\vec{\imath},\vec{\jmath},\vec{k}\)}
\providecommand{\Ouv}{\(\mathrm{O};\vec{u},\vec{v}\)}
\providecommand{\intervFF}[2]{\ensuremath{[#1;#2]}}
\renewcommand{\euro}{\ifmmode\text{\texteuro}\else\texteuro\fi}
\DeclareUnicodeCharacter{00B0}{\ifmmode^{\circ}\else\textdegree\fi}
\DeclareUnicodeCharacter{20AC}{\ifmmode\text{\texteuro}\else\texteuro\fi}
% Énoncés : Éducation nationale. Transcription et corrigés : APMEP, Denis Vergès et contributeurs. Découpage et classement : CoopMaths, Sébastien Lozano, Éric Elter et contributeurs. Adaptation de la présentation : Maths Explorer.
% CC BY-SA : https://coopmaths.fr/www/about/
% https://coopmaths.fr/alea/static/bac/2026/tex/bac_2026_06_sujet1_asie_2_cor.tex
\begin{document}
\noindent\textit{APMEP / CoopMaths. CC BY-SA. Présentation adaptée par Maths Explorer.}\par\medskip
\begin{enumerate}
\item Affirmation 1 : Vraie

Pour tout $x$ réel strictement supérieur à 1, on a : $x-1>0\Longrightarrow (x-1)=\sqrt{(x-1)^2}$.

On en déduit :
\[
\begin{aligned}
\forall x\in ]1~;+\infty[,\quad f(x)&=\dfrac{x-1}{\sqrt{x^2-1}}=\dfrac{\sqrt{(x-1)^2}}{\sqrt{x^2-1}}=\sqrt{\dfrac{(x-1)^2}{x^2-1}}\\
&=\sqrt{\dfrac{(x-1)(x-1)}{(x-1)(x+1)}}=\sqrt{\dfrac{x-1}{x+1}}=\sqrt{\dfrac{x\left(1-\frac1x\right)}{x\left(1+\frac1x\right)}}=\sqrt{\dfrac{1-\frac1x}{1+\frac1x}}.
\end{aligned}
\]
Par limites usuelles, on a : $\displaystyle\lim_{x\to +\infty}\dfrac1x=0$.

Par limite de la somme on en déduit : $\displaystyle\lim_{x\to +\infty}1-\dfrac1x=1$ et $\displaystyle\lim_{x\to +\infty}1+\dfrac1x=1$.

Par limite du quotient, on a ensuite :
\[
\lim_{x\to +\infty}\dfrac{1-\frac1x}{1+\frac1x}=\dfrac11=1.
\]
Et enfin, par composition :
\[
\lim_{x\to +\infty}f(x)=\lim_{x\to +\infty}\sqrt{\dfrac{1-\frac1x}{1+\frac1x}}=\lim_{y\to 1}\sqrt{y}=\sqrt1=1.
\]
La fonction $f$ admet bien 1 pour limite en $+\infty$.

\item Affirmation 2 : Vraie

Pour tout entier naturel $n$, posons $P_n$ l'affirmation : « $w_n=(n+1)^2$ ».

Démontrons par récurrence que les affirmations $P_n$ sont toutes vraies.

\begin{itemize}
\item Initialisation : On a d'une part : $w_0=1$;

et d'autre part : $(0+1)^2=1$.

On constate donc que l'affirmation $P_0$ est vraie.

\item Hérédité : Soit $n$ un entier naturel. On suppose que l'affirmation $P_n$ est vraie.
\[
\begin{aligned}
P_n \text{ est vraie }&\Longrightarrow w_n=(n+1)^2\\
&\Longrightarrow w_n+2n+3=(n+1)^2+2n+3\\
&\Longrightarrow w_{n+1}=(n+1)^2+2n+3 \quad \text{en utilisant la relation de récurrence de }(w_n)\\
&\Longrightarrow w_{n+1}=n^2+2n+1+2n+3\\
&\Longrightarrow w_{n+1}=n^2+4n+4\\
&\Longrightarrow w_{n+1}=(n+2)^2\\
&\Longrightarrow P_{n+1}\text{ est vraie.}
\end{aligned}
\]

\item Conclusion : On a montré que l'affirmation $P_0$ est vraie, et que pour un $n$ entier naturel, si $P_n$ est vraie, alors $P_{n+1}$ est vraie.
\end{itemize}

En vertu du principe de récurrence, on en déduit que, pour tout entier naturel $n$, l'affirmation est vraie, et donc on a établi une expression explicite du terme général de la suite $(w_n)$ :
\[
\forall n\in\N,\quad w_n=(n+1)^2.
\]

\item Affirmation 3 : Vraie

Soit $p$ un nombre réel tel que $0<p<1$, et $X$ une variable aléatoire qui suit la loi binomiale de paramètres 3 et $p$.

On a, par propriété :
\[
\begin{aligned}
P(X=1)&=\binom31\times p^1\times (1-p)^{3-1}\\
&=3\times p\times (1-p)^2\\
&=3p(1-2p+p^2)\\
&=3p-6p^2+3p^3.
\end{aligned}
\]
La probabilité de l'évènement $(X=1)$ est bien : $P(X=1)=3p-6p^2+3p^3$.

\item Affirmation 4 : Fausse

Soit $n$ un entier naturel $n$ non nul.

Une primitive sur $[0~;~1]$ de $x\longmapsto \e^{nx}$ est : $x\longmapsto \dfrac1n\e^{nx}$.

On a donc :
\[
v_n=\int_0^1 \e^{nx}\,\mathrm{d}x=\left[\dfrac1n\e^{nx}\right]_0^1=\dfrac1n\e^{n\times 1}-\dfrac1n\e^{n\times 0}=\dfrac{\e^n-1}{n}\neq \dfrac{\e^n}{n}.
\]

\item Affirmation 5 : Fausse

Chacune des 16 cases est coloriée, donc à condition de numéroter les 16 cases, en associant à chaque case la couleur utilisée, cela revient à associer aux entiers de 1 à 16 une couleur.

Chaque coloriage est ainsi assimilable à une 16-liste d'éléments choisis dans l'ensemble des couleurs, qui est de cardinal 3 (puisque trois couleurs sont à notre disposition). Les éléments peuvent se répéter, puisque nous devons colorier 16 cases avec trois couleurs, donc on devra utiliser la même couleur plusieurs fois.

Le nombre de $k$-listes d'éléments d'un ensemble de cardinal $n$ est : $n^k$.

Il y a donc $3^{16}=43\,046\,721$ coloriages différents possibles.

$\binom{16}{3}=560\neq 43\,046\,721$. Ce nombre pourrait correspondre au nombres de façons qu'il y a de choisir 3 cases parmi les 16 cases du quadrillage.

Remarque : On peut aussi raisonner en disant qu'on a 3 choix de couleur pour la première case, à multiplier par 3 choix pour la deuxième case, et ainsi de suite, on retrouve bien $3^{16}$ choix possibles pour colorier nos 16 cases.
\end{enumerate}

\end{document}
